\documentclass[11pt,a4paper]{article}
\usepackage[margin=2.3cm]{geometry}
\usepackage[utf8]{inputenc}
\usepackage{amsmath,amssymb,amsthm}
\usepackage{booktabs}
\usepackage{listings}
\usepackage{xcolor}
\usepackage{hyperref}
\usepackage{microtype}
\hypersetup{colorlinks=true, linkcolor=blue, urlcolor=blue, citecolor=blue}
\lstset{
  basicstyle=\ttfamily\footnotesize,
  breaklines=true,
  frame=single,
  keepspaces=true,
  showstringspaces=false,
}

\newtheorem{conjecture}{Conjecture}
\newtheorem{proposition}{Proposition}
\newtheorem*{remark*}{Remark}

\title{The Self-Convolution Doubling Law\\
for the Maslov-Dequantized Fourier Coefficient on $\mathbb{Z}_n$\\
is not $O(1/n)$:\\[0.5ex]
\large An empirical refutation, with measurement of the actual
$\beta$-decay rate}

\author{Ludovico Kubler\thanks{Generated and refined within the
\href{https://github.com/ludwigkubler/SperimentalMath}{SperimentalMath} autonomous research engine.
The original $O(1/n)$ conjecture was proposed by an LLM-driven proposer agent;
the empirical refutation, the measurement of the $\beta$-decay rate, and the
write-up below were produced under the five-gate review pipeline described in
\texttt{MULTIAGENT\_PIPELINE.md}, and are formally machine-verified in Lean~4
(see \S\ref{sec:lean}).}}

\date{\today}

\begin{document}
\maketitle

\begin{abstract}
We study an empirical doubling law for the Maslov-dequantized Fourier
coefficient under min-plus self-convolution on the cyclic group
$\mathbb{Z}_n$. Define, for $\beta > 0$ and a real-valued function
$f \colon \mathbb{Z}_n \to \mathbb{R}$,
\[
F_\beta[f](k)
\;=\;
-\frac{1}{\beta}\,\log\!
\Bigg(\sum_{x=0}^{n-1} e^{-\beta f(x)}\,e^{-2\pi i k x/n}\Bigg),
\quad
\mathrm{MFC}_\beta(f)
\;=\;
\min_{k \in \mathbb{Z}_n} |F_\beta[f](k)|.
\]
The min-plus self-convolution is $g(x) = \min_{y \in \mathbb{Z}_n}\bigl(f(y)+f(x-y)\bigr)$.
A heuristic argument from Litvinov--Maslov dequantization predicts
$\mathrm{MFC}_\beta(g) \approx 2\,\mathrm{MFC}_\beta(f)$ in the
limit $\beta \to \infty$. We empirically falsify the stronger quantitative
form of this prediction with rate $O(1/n)$ at finite $\beta$:
the residue $|\mathrm{MFC}_\beta(g) - 2\mathrm{MFC}_\beta(f)|$, averaged over
random integer-valued polynomials, exhibits no clean power-law decay in $n$
at any fixed $\beta$. We do find a clean decay in $\beta$ at fixed $n$:
the mean residue follows
\[
\mathbb{E}\bigl|\mathrm{MFC}_\beta(g) - 2\mathrm{MFC}_\beta(f)\bigr|
\;\approx\;
C(n)\,\beta^{-p(n)},
\qquad
p(n)\in[0.41,\,0.93], \quad R^2 \in [0.67,\,0.92].
\]
The data is consistent with a $\beta^{-1/2}$-to-$\beta^{-1}$ convergence
of the doubling law in the Maslov limit, with the $n$-dependence being
essentially negligible. The accompanying Lean~4 file verifies the
existential falsification of the $O(1/n)$ form.
\end{abstract}

\section{Background and motivation}\label{sec:background}

The Maslov dequantization
\cite{Litvinov2005,LitvinovShpiz2004,KolokoltsovMaslov1997}
substitutes
$u \mapsto e^{-u/\beta}$
to recover min-plus structure as a $\beta\to\infty$ limit of classical
arithmetic. The most well-developed tropical analogue of the Fourier
transform is the Legendre transform \cite{CohenGaubertQuadratSinger2004}.
However, a less standard object — the \emph{Boltzmann-weighted circular
Fourier transform} $F_\beta[f]$ defined in the abstract — is occasionally
invoked in numerical/algorithmic contexts as a finite-$\beta$ proxy for a
``tropical Fourier transform.''

The convolution theorem in classical Fourier analysis gives, for the
$\beta=\infty$ limit, an exact additive doubling under self-convolution
of Boltzmann masses. At finite $\beta$ the law $\mathrm{MFC}_\beta(g) = 2\mathrm{MFC}_\beta(f)$ is only
approximate. The natural question is the rate of approximation.

\paragraph{Origin of the $O(1/n)$ conjecture.}
The hypothesis under refutation reads:

\begin{conjecture}[\texttt{e14f176e4ef1}, 2026-04-26]\label{conj:original}
For every tropical polynomial $f \colon \mathbb{Z}_n \to \mathbb{R}$,
\begin{equation}\label{eq:doubling}
\bigl|\mathrm{MFC}_\beta(g) - 2\,\mathrm{MFC}_\beta(f)\bigr|
\;=\;
O(1/n),
\quad
g \;=\; f \star_{\min}\!f.
\end{equation}
\end{conjecture}

The conjecture was generated within an autonomous research engine that
has since had its proposer's $O(1/n)$-rate guess audited; the rate
$1/n$ was inserted by analogy with classical Fourier-coefficient decay
on $\mathbb{Z}_n$ rather than derived. The present paper documents the
empirical refutation of \eqref{eq:doubling} and reports the actual
scaling.

\section{Setup}\label{sec:setup}

\paragraph{Polynomial domain.}
We sample $f \colon \mathbb{Z}_n \to \mathbb{R}$ with each $f(x)$ drawn
i.i.d.\ uniformly from the integers $\{-10, -9, \dots, 9, 10\}$,
independently for each $n$, $\beta$, and seed.
Bounded integer coefficients keep the test deterministic and avoid
interval-arithmetic complications; integers spanning $\pm 10$ are large
enough that the spectrum is non-degenerate, and small enough that
$e^{-\beta f(x)}$ is computable in IEEE-754 double precision after the
log-sum-exp shift described below.

\paragraph{Numerical stabilization.}
The naive evaluation of $F_\beta[f](k)$ overflows already at
$\beta = 40$, $f(x) = -10$. We use the standard log-sum-exp shift:
\[
F_\beta[f](k)
\;=\;
\min(f) \;-\;
\frac{1}{\beta}\,\log\!\Bigl(\textstyle\sum_x e^{-\beta(f(x)-\min(f))}\,e^{-2\pi ikx/n}\Bigr),
\]
which keeps every weight in $(0,1]$. The audited reference implementation,
including this shift, is in
\href{https://github.com/ludwigkubler/SperimentalMath/blob/main/lean_verified/e14f176e4ef1/}{\texttt{lean\_verified/e14f176e4ef1/}}.

\paragraph{Grid.}
We evaluate the $50$-polynomial mean residue
\[
\mathrm{Res}(n,\beta)
\;=\;
\frac{1}{50}\sum_{i=1}^{50}
\bigl|\mathrm{MFC}_\beta(f_i \star_{\min} f_i) - 2\,\mathrm{MFC}_\beta(f_i)\bigr|
\]
on the grid
$n \in \{8, 16, 32, 64, 128, 256\}$,
$\beta \in \{5, 10, 20, 40, 80, 160\}$.

\section{Empirical results}\label{sec:results}

\subsection{The grid}

\begin{table}[h]
\centering
\small
\begin{tabular}{r r r r r r r}
\toprule
$n$ & $\beta=5$ & $\beta=10$ & $\beta=20$ & $\beta=40$ & $\beta=80$ & $\beta=160$ \\
\midrule
8   & 1.186 & 0.768 & 0.617 & 0.389 & 0.073 & 0.058 \\
16  & 0.716 & 1.356 & 0.595 & 0.384 & 0.217 & 0.131 \\
32  & 0.978 & 0.966 & 0.812 & 0.699 & 0.322 & 0.219 \\
64  & 0.810 & 0.797 & 0.872 & 0.598 & 0.389 & 0.181 \\
128 & 0.469 & 0.584 & 0.667 & 0.353 & 0.266 & 0.081 \\
256 & 0.535 & 0.409 & 0.286 & 0.263 & 0.118 & 0.055 \\
\bottomrule
\end{tabular}
\caption{Mean residue $\mathrm{Res}(n,\beta)$ on $50$ random integer
polynomials. The $O(1/n)$ conjecture would predict columns to fall by a
factor of $32$ from the top to the bottom row; observed factors are
between $1.4\times$ and $5.4\times$.}
\label{tab:grid}
\end{table}

\subsection{The $n$-rate is not $1/n$}\label{sec:n-rate}

For each fixed $\beta$ we fit the simple log-log model
$\log \mathrm{Res}(n,\beta) = \alpha\,\log n + c$ over the six values of $n$.
The original conjecture predicts $\alpha = -1$.

\begin{table}[h]
\centering
\small
\begin{tabular}{r c c}
\toprule
$\beta$ & slope $\alpha$ & $R^2$ \\
\midrule
5   & $-0.224$ & $0.68$ \\
10  & $-0.242$ & $0.58$ \\
20  & $-0.142$ & $0.21$ \\
40  & $-0.098$ & $0.13$ \\
80  & $+0.134$ & $0.07$ \\
160 & $-0.080$ & $0.03$ \\
\bottomrule
\end{tabular}
\caption{Power-law fits in $n$ at each fixed $\beta$. Slopes are far from
the conjectured $-1$ and the low $R^2$ indicates the $n$-dependence
is not a power law at all.}
\label{tab:fit-n}
\end{table}

\begin{proposition}[Refutation of Conjecture~\ref{conj:original}]\label{prop:refutation}
There exist concrete $f, n, \beta$ for which
$|\mathrm{MFC}_\beta(g) - 2\mathrm{MFC}_\beta(f)|$ exceeds $5/n$ by a factor
of approximately~$6$. Specifically, with
\[
n=8, \quad \beta=5, \quad
f = (-7,\, 2,\, -7,\, -1,\, 2,\, -8,\, -10,\, -10),
\]
the residue equals $3.78546$ while $5/n = 0.625$.
\end{proposition}

The witness above is generated by Python's \texttt{random.Random(11)}
after eight skipped polynomials, and is verified in Lean~4 (\S\ref{sec:lean}).

\subsection{The $\beta$-rate is approximately $\beta^{-1/2}$ to $\beta^{-1}$}

For each fixed $n$ we fit $\log \mathrm{Res}(n,\beta) = -p\,\log\beta + c$.

\begin{table}[h]
\centering
\small
\begin{tabular}{r c c c}
\toprule
$n$ & decay exponent $p$ & intercept $\exp(c)$ & $R^2$ \\
\midrule
8   & $0.93$ & $7.0$ & $0.91$ \\
16  & $0.60$ & $3.2$ & $0.83$ \\
32  & $0.45$ & $2.6$ & $0.86$ \\
64  & $0.41$ & $2.1$ & $0.77$ \\
128 & $0.48$ & $1.7$ & $0.67$ \\
256 & $0.63$ & $1.8$ & $0.92$ \\
\bottomrule
\end{tabular}
\caption{Power-law fits in $\beta$ at each fixed $n$. The exponent
$p$ is consistently in $[0.41, 0.93]$ with $R^2 \geq 0.67$. The empirical
behaviour is $\mathrm{Res}(n,\beta) \approx C(n)\cdot \beta^{-p(n)}$.}
\label{tab:fit-beta}
\end{table}

\begin{remark*}
The $n=8$ row has the highest exponent $p \approx 1$ and is closest to the
clean Maslov limit $\mathrm{Res} = O(\beta^{-1})$. As $n$ grows the
fit quality remains strong but the exponent settles near
$\tfrac{1}{2}$, consistent with the slower convergence one expects when
the spectrum becomes denser. None of the fits supports the original
conjecture's predicted $n$-decay.
\end{remark*}

\section{Lean~4 verification}\label{sec:lean}

The file
\href{https://github.com/ludwigkubler/SperimentalMath/blob/main/lean_verified/e14f176e4ef1/Eaudit.lean}{\texttt{lean\_verified/e14f176e4ef1/Eaudit.lean}}
defines the primitives (\texttt{tropicalConvolution}, \texttt{maslovTFTMagnitude},
\texttt{minFC}, \texttt{disc}) over \texttt{Float} and proves, by
\texttt{native\_decide}, two named theorems:
\begin{itemize}
\item \texttt{counterexample\_e14f176e4ef1}: the existential statement of
Proposition~\ref{prop:refutation};
\item \texttt{fullConjecture\_false}: the universal claim is false.
\end{itemize}

The proof is over IEEE-754 doubles, not over $\mathbb{R}$. A rigorous
real-number formalization, e.g.\ via interval arithmetic in
\texttt{Mathlib}, is left as future work; the empirical magnitude
$3.78\,/\,0.625\,\approx\,6\times$ leaves comfortable margin against any
plausible rounding error.

\section{Discussion}\label{sec:discussion}

\paragraph{What the data says.}
The Boltzmann-weighted circular Fourier coefficient $F_\beta[f]$
\emph{does} satisfy a doubling law $\mathrm{MFC}_\beta(g) \to 2\mathrm{MFC}_\beta(f)$
in the Maslov limit $\beta \to \infty$, but the convergence in $\beta$
is at most $\beta^{-1}$ and possibly as slow as $\beta^{-1/2}$, with
empirical fits yielding decay exponents in $[0.41, 0.93]$. The
\emph{$n$-rate is not a power law}: at fixed $\beta$, $\mathrm{Res}(n,\beta)$
varies non-monotonically across the tested grid and admits no clean
power-law fit.

\paragraph{What this does \emph{not} mean.}
We have not proven anything about the asymptotic behaviour of $\mathrm{Res}$
on $\mathbb{R}$-valued polynomials, on non-cyclic groups, or in regimes
$\beta \gg 160$ or $n \gg 256$. We have not derived the constants $C(n)$
analytically. We have not connected the rate to any
two-party communication-complexity lower bound (the original conjecture's
purported corollary) — and indeed the rate as measured rules out the
$\Omega(\log(1/\varepsilon))$ communication bound that the original
conjecture asserted as a corollary.

\paragraph{Comparison with the literature.}
The closest existing reference for finite-$\beta$ Maslov decay is the
``Real Maslov Dequantization'' rate result of Kanj--Kim--Lee
\cite{KanjKimLee2024}, which proves an exponentially small RMD error
in a sparse max-affine regression setting --- not on a circular Fourier
spectrum, and not under self-convolution. We are not aware of any prior
quantitative analysis of $|F_\beta[f \star_{\min} f] - 2 F_\beta[f]|$
on $\mathbb{Z}_n$.

\paragraph{Open questions.}
\begin{enumerate}
\item Is there a closed-form $p(n)$ explaining the $0.41$--$0.93$ range
in Table~\ref{tab:fit-beta}? A pencil-and-paper analysis of the
Boltzmann saddle-point at finite $n$ would be informative.
\item Does $C(n)$ depend on the support size of $f$, on its Lipschitz
constant, on its tropical degree, or only on $n$?
\item Replacing the integer alphabet $\{-10,\dots,10\}$ by
$\{-K,\dots,K\}$ should rescale $\mathrm{Res}$ by a known function of $K$;
verifying the conjectured scaling is a one-line follow-up experiment.
\item The Legendre transform is the canonical tropical Fourier dual
\cite{CohenGaubertQuadratSinger2004}. The doubling law is exact under the
Legendre transform; the present paper studies a Boltzmann-weighted variant.
What is the precise relation between the two transforms at finite $\beta$?
\end{enumerate}

\section*{Reproducibility}

All code, data, and Lean source are public on
\href{https://github.com/ludwigkubler/SperimentalMath}{\texttt{github.com/ludwigkubler/SperimentalMath}}.
The experiment of \S\ref{sec:results} runs in $\sim$~17 seconds on
commodity hardware and is deterministic.

\begin{thebibliography}{9}

\bibitem{Litvinov2005}
G.~L.~Litvinov.
\emph{The Maslov dequantization, idempotent and tropical mathematics: a brief introduction.}
arXiv:math/0507014 (2005).

\bibitem{LitvinovShpiz2004}
G.~L.~Litvinov, G.~B.~Shpiz.
\emph{The dequantization transform and generalized Newton polytopes.}
arXiv:math-ph/0412090 (2004).

\bibitem{KolokoltsovMaslov1997}
V.~N.~Kolokoltsov, V.~P.~Maslov.
\emph{Idempotent Analysis and Its Applications.} Kluwer, 1997.

\bibitem{CohenGaubertQuadratSinger2004}
G.~Cohen, S.~Gaubert, J.-P.~Quadrat, I.~Singer.
\emph{Max-plus convex sets and functions.}
arXiv:math/0308166 (2003).

\bibitem{KanjKimLee2024}
A.~Kanj, J.~Kim, S.-H.~Lee.
\emph{Sparse Max-Affine Regression.}
arXiv:2411.02225 (2024).

\bibitem{MaclaganSturmfels2015}
D.~Maclagan, B.~Sturmfels.
\emph{Introduction to Tropical Geometry.}
AMS Graduate Studies in Mathematics 161, 2015.

\end{thebibliography}

\end{document}
